% Claim template. Copy the pattern, not the numbers.
%
% 1. The number lives in a generated include, never in this file:
%      \input{generated/tables/<name>.tex}
% 2. CLAIMS.json holds the traceability entry:
%      {
%        "id": "C-007",
%        "value": "0.31",
%        "tex": "paper/sections/results.tex",
%        "artifact": "fixed_case_studies/stocks/results/<run>/manifest_main__2023-01-01__seed42.json",
%        "artifact_sha256": "<sha256 of the manifest>",
%        "commit": "<commit that produced it>",
%        "status": "verified"
%      }
% 3. Prose cites the claim:
%
%      Vangja reaches a median relative MAE of 0.31~[C-007] over the
%      confirmation origins (24 target$\times$origin units, persistence
%      denominator above the $\varepsilon=10^{-3}$ floor, scaled target),
%      against a frozen configuration -- no test-horizon selection.
%
% 4. Check: check_claims.py --paper paper/paper.tex --claims CLAIMS.json --gate

\begin{table}[htbp]
    \centering
    \input{generated/tables/stocks_relmae.tex}
    \caption{Relative MAE against persistence (median over target$\times$origin
    units). Values below 1 are better than persistence. Confirmation origins;
    frozen configuration; includes are machine-generated from
    \texttt{manifest\_*.json} at commit \texttt{[C-007.commit]}.}
    \label{tab:stocks_relmae}
\end{table}
